%! Author = lili
%! Date = 4/27/26

% beim erwartungswert trägt 0 nichts bei, bei der varianz schon? 2026 Folie


\chapter{Erwartungswert und Varianz von Zufallsvariablen}\label{ch:erwartungswert-und-varianz-von-zufallsvariablen}
\untit{Vorlesung 30. Mai 2025 (V07 2025)und 05. Juni 2026 (V07 2026)}


\section*{Lernziele}

\subsection*{Konzepte}
\begin{itemize}
    \item Erwartungswert einer Zufallsvariablen (gewichteter Mittelwert)
    \item Varianz als weitere Kennzahl
\end{itemize}

\subsection*{Fakten}
\begin{itemize}
    \item Zwei verschiedene Darstellungen des Erwartungswerts.
    \item Der Erwartungswert ist linear.
\end{itemize}

\subsection*{Wissen}
\begin{itemize}
    \item Existenz des Erwartungswerts prüfen können.
    \item Spezialfall: \(\zvx (\eraum)\)ist eine endliche Menge (insbesondere erfüllt, wenn \(\eraum\)endlich ist).
    \item Spezialfall: \(\zvx  \geq 0\).
    \item Erwartungswerte berechnen können.
    \item Erwartungswert und Varianz für die wichtigsten Verteilungen kennen
    \item Erwartungswert von Produkten zum Widerlegen von Unabhängigkeit
    \item Erwartungswert von Produkten bei Unabhängigkeit berechnen können
    \item Markov-Ungleichung kennen und anwenden können
\end{itemize}

%%%%%%%%%%%%%%%%%%%%%%%
% Grundgedanke
%%%%%%%%%%%%%%%%%%%%%%%

\section*{Grundgedanke} \txc{(2025)}
Die Verteilung charakterisiert eine \Gls{zv} $\zvx$ eindeutig (Im Sinne ihrer statistischen Eigenschaften: Wir können Wahrscheinlichkeiten
berechnen wie z.B. $\pe(\zvx  \geq 2)$, und diese Wahrscheinlichkeiten hängen nur von der
Verteilung ab, nicht von der Wahl des Wahrscheinlichkeitsraums $\wraum$.). \\
Sind $\zvx$ und $\zvy$ zwei \Glspl{zv} \un{mit derselben Verteilung}, so gilt
\[
    \pe(\zvx  \in A)= \pex(A)= \pey(A)= \pe(\zvy \in A)
\]
für jedes $A$. \\
Es ist nicht immer klar, welche Verteilung zu einem sinnvollen Modell führt.
Auch sind die Parameter der Verteilung nicht unbedingt bekannt.
Siehe \txc{(2026)}
\begin{itemize}
    \item \textbf{Binomialverteilung:}
    \[
        \Binomv(k)
        =
        \binom{n}{k}
        p^k (1-p)^{n-k},
        \qquad
        k \in \{0,\dots,n\}.
    \]

    \item \textbf{Geometrische Verteilung:}
    \[
        G_p(n)
        =
        p(1-p)^{n-1},
        \qquad
        n \in \nz.
    \]
    % TODO gloss entry für geom verteilung

    \item \textbf{Poisson-Verteilung:}
    \[
        P_{\lambda}(k)
        =
        \frac{\lambda^k}{k!} e^{-\lambda},
        \qquad
        k \in \nz_0.
    \]

    \item \textbf{Exponentialverteilung:}
    \[
        f(x)
        =
        \ele e^{-\lambda x}\,
        \chaf_{[0,\infty)}(x),
        \qquad
        x \in \rz.
    \]
\end{itemize}
Wie können wir trotzdem Vorhersagen treffen? \\ \newline
Dies führt uns in die Statistik (später).
% F4

\paragraph{Typische Fragestellungen:} Wie \dots \textit{im Mittel}? \txc{(2025)} \\
\begin{itemize}
    \item Wie groß ist der mittlere Gewinn bei einem Glücksspiel?
    \item Wie lange muss man bei einem wiederholten Zufallsexperiment im Mittel
    warten, bis ein bestimmtes Ereignis eintritt?
    \item Wie viele weitere Personen steckt eine infizierte Person im Mittel an?
\end{itemize}

\begin{bsp}
    \textbf{Ansteckungen} \txc{(2026)}
    \begin{itemize}
        \item Die \textbf{Basisreproduktionszahl} \(R_0\)gibt an, wie viele Personen
        eine infizierte Person im Mittel ansteckt, wenn es in der Bevölkerung keine
        Immunität gibt.
        Gemäß RKI gilt:
        \[
            R_0 \in [2.8,\,3.8]
        \]
        (vor möglicherweise ansteckenderen Mutationen).

        \item Die \textbf{zeitabhängige Reproduktionszahl}
        \[
            R = R(t)
        \]
        berücksichtigt Effekte durch Hygienemaßnahmen, Impfungen und Immunität.
        Sie gibt an, wie viele Personen eine infizierte Person in einem Zeitfenster
        gegebener Länge im Mittel ansteckt.

        \item Die \textbf{Inzidenz} setzt Neuinfektionen in Relation zur
        Bevölkerungsgröße.
        Die \emph{7-Tage-Inzidenz} gibt beispielsweise an,
        wie viele neu infizierte Personen es in den vergangenen sieben Tagen
        pro \(100\,000\)Einwohner gegeben hat.
    \end{itemize}
\end{bsp}

\begin{bsp}
    \textbf{1 $\times$ Würfeln mit einem fairen Würfel} \txc{(2025)} \\
    \txc{Da herrscht eine Gleichverteilung, daher ist die erwartete Augenzahl ist das arithmetische Mittel.} \\
    \[
        \eraum = \{1,\dots,6\}
    \]
    mit Gleichverteilung \(\pe\).

    Die Zufallsvariable
    \[
         \zvx: \eraum \to \rz,
        \qquad
        \zvx(\ele)=\ele,
    \]
    besitzt das Bild
    \[
        \zvx(\eraum)=\{1,\dots,6\}\subseteq\rz.
    \]

    Für einen fairen Würfel ist der Erwartungswert die mittlere Augenzahl, also das
    arithmetische Mittel:
    \[
        \E{\zvx }
        =
        \frac{1+2+3+4+5+6}{6}
        =
        3.5.
    \]

    Die erwartete Augenzahl entspricht der mittleren Zahl von Schritten, die man
    in einem Würfelspiel pro Wurf vorwärtsgehen würde.
\end{bsp}

\begin{bsp}
    \textbf{Spiel, bei dem nicht alle Ausgänge gleichwahrscheinlich sind} \txc{(2025)} \\
    1 $\times$ Würfeln eines fairen Würfels \\
    Auszahlung:
    \begin{itemize}
        \item $k$ Euro, falls 2 gewürfelt
        \item 6 Euro, falls $2, \cdots , 6$ gewürfelt
        \item \(
        \eraum = \{1,\dots,6\}
        \)mit Gleichverteilung \(\pe\).
    \end{itemize}
    \glssymbol{eraum} = $\kl{1, \dots, 6}$ mit Gleichverteilung \pe  \\
    Die Auszahlung wird beschrieben durch die \Gls{zv} $\zvy$ mit
    \[
        \zvy(\ele)= k \cdot \chaf_{\kl{1}}(\ele)+ 6 \cdot \chaf_{\kl{2, \dots, 6}}(\ele)\quad \faoio
    \]
    Erwartete Auszahlung als arithmetisches Mittel:
    \mab{\E{\zvy} = \sum_{\oio} \zvy(\ele)\glssymbol{wfk}(\ele)= \frac{k+ 5 \cdot 6}{6} = 6 \cdot \frac{5}{6} + k \cdot \frac{1}{6}  = 6 \cdot \pe(\zvy=6)+ k \cdot \pe(\zvy=k)= \sum_{I\in \zvy(\glssymbol{eraum})} I \cdot \pe(\zvy=I)}

\end{bsp}

\karkar{bem_verteilung_legt_stat_eigenschaften_fest}

%%%%%%%%%%%%%%%%%%%%%%%
% Definition und Existenz des Erwartungswerts
%%%%%%%%%%%%%%%%%%%%%%%

\section{Definition und Existenz des Erwartungswerts}

\begin{defi} \txc{(2025)}
    Der \Gls{ew} einer \zv $\zvx$ mit abzählbarem Bildbereich $\zvx(\eraum)\in \rz$ ist gegeben durch
    \[
        \E{\zvx } = \mathbb{E} \zvx = \sum_{x\in \zvx(\eraum)} x \cdot \pe (\zvx =x)
    \]
    insofern die Reihe \glsverb{abskonv}, also \ma{ \sum_{x\in \zvx(\eraum)} |x| \cdot \pe (\zvx =x)< \infty}. \\
    Wir sprechen in diesem Fall von der Existenz des \Gls{ew}s. \\ \newline
    Der \Gls{ew} einer \zv $\zvx$ mit einer \dichte \glssymbol{dichte} ist gegeben durch
    \mab{\E \zvx = \int_{- \infty}^{\infty} x \glssymbol{dichte}(x)\dx = \lim_{N \rightarrow \infty} \int_{-N}^{N} x \glssymbol{dichte}(x)\dx }
    insofern \ma{\int_{- \infty}^{\infty} |x| \glssymbol{dichte}(x)\dx < \infty}. \\
    Wir sprechen in diesem Fall von der Existenz des Erwartungswerts.

    \paragraph{Hinweis:} Für \glspl{dichte} \dcht, die gerade Funktionen sind, gilt \(\int_{-N}^N x \dcht (x)\dx = 0\)für
    alle N, auch wenn der Erwartungswert ggf.
    gar nicht existiert.
\end{defi}

\begin{bem} \txc{(2026)}
    Absolute Konvergenz der Reihe, die den Erwartungswert definiert, garantiert,
    dass wir die Summanden beliebig umsortieren dürfen.
    Es gilt daher
    \[
        \txo{\E{\zvx } \overset{Def.}{=}} \sum_{x\in \zvx(\eraum)} x \pe(\zvx =x)= \sum_{x\in \eraum} \zvx(\ele)\wfk(\ele)
    \]
    sofern eine der Reihen absolut konvergent ist.
    (Vgl.
    obiges Bsp.
    mit dem Spiel.)

    Wenn $\eraum$ oder $\zvx(\eraum)$ endlich ist, so existiert der Erwartungswert automatisch.

    Nimmt $\zvx$ nur nichtnegative Werte an, können wir einfach anfangen, den
    Erwartungswert auszurechnen, weil
    \[
        \sumxinXO  |x| \pe(\zvx =x)= \sumxinXO x \pe(\zvx =x)
    \]
    bzw.
    \[
        \intinfs |x| \dcht (x)\dx = \intinfs x \dcht (x)\dx \txc{ = \int^{\infty}_0 x \dcht (x)\dx }
    \]
\end{bem}

\karkar{def_erwartungswert_diskret}
\karkar{def_erwartungswert_dichte}
\karkar{bem_gerade_dichte_fallstrick}
\karkar{bem_umordnung_und_summationsbereich_ew}
\karkar{bem_existenz_ew_spezialfaelle}

%%%%%%%%%%%%%%%%%%%%%%%
% Erwartungswerte der wichtigsten Verteilungen
%%%%%%%%%%%%%%%%%%%%%%%

%\subsection{Erwartungswerte der wichtigsten Verteilungen}

\begin{bsp}
    \textbf{Erwartungswert der Poisson-Verteilung} \txc{(2026)}  \\
    $\zvx$ Poisson-verteilt, d.h., es existiert ein $\lambda > 0$ mit
    \[
        \wfk_{\zvx}(k)= \pe (\zvx =k)= \frac{\lambda^k}{k!} e^{-\lambda} \quad \text{für alle } k \in \nz_0 = \{0, 1, 2, 3, \dots \}
    \]
    Der Erwartungswert existiert, weil die folgende Rechnung ein endliches Ergebnis hat: \\
    (Wir verwenden in dieser Argumentation, dass $\zvx$ nur nichtnegative Werte hat, d.h., $\zvx \geq 0$.)
    \begin{align*}
        \E{\zvx } & \overset{\txo{*1}}{=} \sum^{\infty}_{k=0} k \pe(\zvx =k)\\
        & \overset{\txo{*2}}{=}  \sum^{\infty}_{k=0} \frac{\lambda^k}{k!} e^{-\lambda} \\
        & = \sum^{\infty}_{\txr{k=1}} \txg{k} \cdot \frac{\txb{\lambda^k}}{\txg{k!}} e^{-\lambda} \\
        & = \txb{\lambda}   \sum^{\infty}_{\txr{k=1}}  \frac{\txb{\lambda^{k-1}}}{\txg{(k-1)!}} e^{-\lambda} \\
        & = \lambda e^{-\lambda} \sum^{\infty}_{\txr{k=1}}  \frac{\lambda^{k-1}}{(k-1)!} \\
        & = \lambda e^{-\lambda}
        \oub{
            \sum_{l=0}^{\infty}
            \frac{\lambda^{\txr{l}}}{\txr{l}!}
        }{=e^{\lambda}} \\
        & = \lambda
    \end{align*}
    \txo{
        \begin{align*}
            *1 & \quad \zvx(\eraum)= \nz_0 \\
            *2 & \quad \pe(\zvx =k)= \wfk_{\zvx}(k)= \dots
        \end{align*}
    }
\end{bsp}

\begin{bsp}
    \textbf{Erwartungswert der geometrischen Verteilung} \txc{(2026)}  \\
    Wie lange müssen wir im Mittel warten, bis der erste Erfolg eintritt?
    (Zum Beispiel bis zur ersten Sechs beim Werfen eines fairen Würfels?
    – Raten Sie!)\\
    $\zvx$ geometrisch verteilt mit Erfolgswahrscheinlichkeit \(p \in (0,1)\), d.h.,
    \[
        \wfk_{\zvx}(k)= \pe (\zvx =k)= p(1-p)^{k-1} \quad \text{für alle } k \in \nz_0 = \{0, 1, 2, 3, \dots \}
    \]
    Der Erwartungswert existiert, weil die folgende Rechnung ein endliches Ergebnis hat (Wir verwenden wieder, dass $\zvx \geq 0$)).
    \txo{
        \[
            h(x)= x^k \Rightarrow h'(x)= k x^{k-1}
        \]
    }
    \begin{align*}
        \E{\zvx }&  \overset{\txo{*1}}{=}   \sum^{\infty}_{k=1} \txr{k} p(1-p)^{\txr{k-1}} \\
        & =  p  \sum^{\infty}_{k=1} \big(\frac{d}{dx} x^{\txr{k}} \big)\big|_{x=1-p} \\
        &  \overset{\txo{*2}}{=}  p  \frac{d}{dx}   \sum^{\infty}_{\txg{k=1}} \big(x^{\txr{k}} \big)\big|_{x=1-p} \\
        & = p  \frac{d}{dx}   \sum^{\infty}_{\txg{k=0}} \big(x^{k} \big)\big|_{x=1-p} \\
        & = p  \frac{d}{dx}  \big(\frac{1}{1-x} \big)\big|_{x=1-p} \\
        & = p  \frac{1}{(1-x)^2}  \big|_{x=1-p} \\
        & = \frac{1}{p}
    \end{align*}
    \txo{
        \begin{align*}
            *1 & \quad \text{Definition Erwartungswert} \\
            *2 & \quad \text{gilt, weil gliedweises Di”erenzieren einer Potenzreihe im Innern ihres Konvergenzradius
            zulässig ist}
        \end{align*}
    }
\end{bsp}

\begin{bsp}
    \textbf{Erwartungswert der uniformen Verteilung auf [a, b]} \txc{(2026)}  \\
    (Raten.
    – Warum ist der Erwartungswert (a + b)/2 ?)\\
    Der Erwartungswert existiert, weil\ldots
    \txc{
        Denn: Zu prüfen ist: \(\intinfs |x| \dcht (x)\dx < \infty  \)für $\dcht(x)= \frac{1}{b-a} \chaf_{[a, b]}(x)$. \\
        Hier ist $\dcht(x)= 0$ außerhalb eines kompakten Intervalls (abgeschlossen und beschränkt).
        \[
            \intinfs |x| \dcht(x)\dx = \intab \frac{|x|}{b-a} \dx \leq \intab \frac{\max\{|a|, |b|\}}{b-a} \dx = \max\{|a|, |b|\} \cdot 1 < \infty
        \]
    }
    \ldots f stetig ist auf dem kompakten Intervall [a, b]
    und f = 0 außerhalb des Intervalls [a, b] gilt.
    \[
        \E{\zvx } = \intab x \frac{1}{b-a} \dx = \frac{1}{2(b-a)} x^2 |^b_{x=a} = \frac{b^2 - a^2}{2(b-a)} = \oub{\frac{a+b}{2}}{*}
    \]
    \txo{* Mittelpunkt des Intervalls} \\
    \txo{
        Erinnerung: $\intab x \dx = \frac{x^2}{2} |^b_{x=a}$
    }
\end{bsp}
\begin{bsp}
    \textbf{Erwartungswert der Exponentialverteilung} \txc{(2026)}  \\
    \txo{
        Erinnerung:}
    \[
        \txo{ F(x)\coloneqq - e^{-\alpha x}, \quad F'(x)\coloneqq \alpha e^{-\alpha x}, \quad \E{\zvx } = \frac{1}{\alpha}}
    \]

    Der Erwartungswert existiert genau dann, wenn die folgende Berechnung einen
    endlichen Wert ergibt (was der Fall ist, wir verwenden wieder, dass $\zvx \geq 0$):
    \begin{align*}
        \E{\zvx } & = \int^{\infty}_0 \underbrace{x}_{\downarrow} \underbrace{\alpha e^{- \alpha x}}_{\uparrow} \dx \\
        & = - \lim_{N \rightarrow \infty} \oub{x e^{- \alpha x}}{=0} |^N_{x=0} +  \int^{\infty}_0 \oub{1 \cdot  e^{- \alpha
        x} }{=f(x)\cdot \frac{1}{\alpha}} \dx \\
        & = 0 + \frac{1}{\alpha} \oub{\intinfs \dcht(x)\dx }{=1} \\
        & = \frac{1}{\alpha}^{\txo{*}}
    \end{align*}
    \txc{
        \begin{align*}
            * \dcht(x)& = \alpha e^{- \alpha x} \chaf_{[0,\infty)}(x)\\
            \Rightarrow \int_0^{\infty} e^{- \alpha x} \dx
            & = \frac{1}{\alpha} \int_0^{\infty} \alpha e^{- \alpha x} \dx \\
            & = \frac{1}{\alpha}  \intinfs \dcht(x)\dx \\
            & = \frac{1}{\alpha}
        \end{align*}
    }
\end{bsp}

%%%%%%%%%%%%%%%%%%%%%%%
% Rechenregeln für den Erwartungswert
%%%%%%%%%%%%%%%%%%%%%%%

\subsection{Rechenregeln für den Erwartungswert}

\begin{satz}
    \textbf{Der Erwartungswert ist linear} \txc{(2026)} \\ % 2026 13
    Geg. \glspl{zv} $\zvx: \eraum \rightarrow \rz $ und  $\zvy: \eraum \rightarrow \rz $auf $\wraum$; Konstanten $a,b \in \rz$. \\
    Annahme: $\E{\zvx }$ und $\E{\zvy}$ existieren. \\
    Dann gelten $a \zvx + b \zvy$ ist ebenfalls eine Zufallsvariable.
    Der Erwartungswert von $a \zvx + b \zvy$ existiert und berechnet sich als
    \[
        \E{a \zvx + b \zvy} = a\E{\zvx } + b\E{\zvy}
    \]

    \paragraph{Beweis (für Zufallsvariable mit abzählbarem Bildbereich)}
    Nehmen wir zunächst an, der Erwartungswert von $a \zvx + b \zvy$ existiert.
    Dann dürfen alle Reihen beliebig umsortiert werden:
    \begin{align*}
        \E{a \zvx + b \zvy} & = \sum_{\oio} [a\zvx(\ele)+ b\zvy(\ele)] \wfk(\ele)\\
        & = a \sum_{\oio} \zvx(\ele)\wfk(ele)+ \sum_{\oio} \zvy(\ele)\wfk(ele)\\
        & = a\E{\zvx } + b\E{\zvy}
    \end{align*}
    Eine analoge Rechnung mit Betragszeichen und Abschätzungen zeigt die
    Existenz des Erwartungswerts:
    \[
        \sum_{\oio} \oub{ | a\zvx(\ele)+ b\zvy(\ele)| }{\leq |a| |\zvx(\ele)| + |b| |\zvy(\ele)|} \wfk(\ele)\leq |a| \oub{\sum_{\oio} |\zvx(\ele)| \wfk(\ele)}{< \infty\ unab.\ Ver.} + |b|  \oub{\sum_{\oio} |\zvy(\ele)| \wfk(\ele)}{< \infty\ unab.\ Ver.} < \infty
    \]
    da nach Voraussetzung $\E{\zvx }$ und $\E{\zvy}$ existieren.
\end{satz}

\begin{bsp}
    \textbf{Augensumme bei mehrfachem Würfeln mit fairem Würfel} \txc{(2026)}
\[
        \eraum = \{1,\dots,6\}^n
    \]
    mit Gleichverteilung \(P\).

    Für \(i=1,\dots,n\)sei
    \[
        \zvx_i(\ele)=\ele_i,
    \]
    wobei \(\ele=(\ele_1,\dots,\ele_n)\in\eraum\).

    Die Zufallsvariable
    \[
        \zvx=\sum_{i=1}^{n} \zvx_i
    \]
    beschreibt die Augensumme \txo{(Und $\zvx_i$ das Ergebnis des i-ten Wurfs)}.

    \textbf{Berechnung der erwarteten Augensumme für zweimaliges Würfeln(\(n=2\)):}

    Direkt über die Verteilung:
    \[
        \E{\zvx }
        =
        \sum_{i=2}^{12}
        i \, \pe(\zvx _1+\zvx_2=i)= \dots
    \]
    was etwas aufwendig ist.

    Alternativ:
    \[
        \begin{aligned}
            \E{\zvx }
            &=
            \sum_{\ele\in\eraum}
            \bigl(\zvx _1(\ele)+\zvx_2(\ele)\bigr)\,\wfk(\ele)\\
            &=
            \frac{1}{36}
            \sum_{\ele_1, \ele_2=1}^{6} [\ele_1 + \ele_2] \\
            &=
            \frac{1}{36}
            \sum_{\ele_1=1}^{6}
            \sum_{\ele_2=1}^{6}
            (\ele_1+\ele_2)\\
            &=
            \frac{1}{36}
            \left[
                6\sum_{\ele_1=1}^{6}\ele_1
            +
            6\sum_{\ele_2=1}^{6}\ele_2
            \right] \\
            &=
            \frac{1}{36}
            \left(
                6\cdot 21 + 6\cdot 21
            \right)
            =
            7.
        \end{aligned}
    \]

    Mit der Linearität des Erwartungswertes erhält man sofort:
    \[
        \E{\zvx }
        =
        \E{\zvx _1}+\E{\zvx _2}
        =
        2\cdot 3,5
        =
        7.
    \]

    Dies verallgemeinert sich unmittelbar auf \(n\)-maliges Würfeln:
    \[
        \E{\zvx }
        =
        \E{\zvx _1+\cdots+\zvx_n}
        =
        \E{\zvx _1} + \cdots+ \E{\zvx _n}
        =
        n\cdot\frac{7}{2}.
    \]
\end{bsp}

\begin{bsp}
    \textbf{Binomialverteilung} \txo{$\E{\zvx } = n \cdot p$} \txc{(2026)}  \\
    \txo{
        $\zvx(\eraum)= \{ 0,\dots, n\}, \quad \# \zvx(\eraum)= n+1 < \infty$}
    \begin{itemize}
        \item \(\zvx \)sei binomialverteilt mit den Parametern \(n\)und \(p\).
        \item Der Erwartungswert$^{\txo{*}}$ von \(\zvx \)existiert, da der Bildbereich von \(\zvx \)
        endlich ist.
    \end{itemize}
    \txo{erwartete Anzahl an Erfolgen in einem n-fach unabhängig wiederholten Bernoulli-Experiment}

    \paragraph{1. Versuch} (mit binomischer Formel und etwas Geschick)
    \[
        \E{\zvx } = \sum^n_{k=0} k \oub{\binom{n}{k} p^k (1-p)^{n-k}}{=\pe(\zvx =x)} = \dots
    \]

    \paragraph{2. Versuch}
    Wähle
    \[
        \eraum=\{0,1\}^n.
    \]

    Die Anzahl der Erfolge lässt sich darstellen als
    \[
        \zvx=\sum_{i=1}^{n} \zvx_i,
    \]
    wobei
    \[
        \zvx_i(\ele)=\ele_i.
    \]

    Für jede Indikatorvariable gilt
    \[
        \E{\zvx _i}
        =
        0\cdot \pe(\zvx _i=0)
        +
        1\cdot \pe(\zvx _i=1)
        =
        p.
    \]

    Mit der Linearität des Erwartungswertes folgt unmittelbar:
    \[
        \E{\zvx }
        =
        \sum_{i=1}^{n} \oub{\E{\zvx _i}}{=p}
        =
        \sum_{i=1}^{n} p
        =
        np.
    \]
\end{bsp}

\begin{bsp}
    \textbf{St. Petersburger Paradoxon (Erwartungswert muss nicht existieren)} \txc{(2026)} \\
    Ein Casino möchte folgendes Spiel anbieten:
    Eine faire Münze wird geworfen, bis zum ersten Mal Kopf erscheint.
    Passiert dies beim n-ten Wurf, so sei die Auszahlung $\zvx = 2^n$
    \[
        \eraum=\nz,
        \qquad
        \wfk(n)=p(1-p)^{n-1},
        \qquad
        p=\frac12.
    \]
    Die Zufallsvariable $\zvx:\eraum\to\rz$ sei definiert durch $ \zvx(n)=2^n$ \txo{(Ausazhlung)} für alle $n\in\eraum$

    \[
        \wfk_{\zvx}(2^n)
        =
        \pe(\oub{\zvx =2^n}{\text{Kopf erstmalig im n-ten Wurf}})
        =
        \wfk(n)
        =
        p(1-p)^{n-1}
        =
        2^{-n},
        \qquad
        n\in \eraum = \nz.
    \]
    \txo{$\wfk(n)= \left(\frac{1}{2}\right)^n = 2^{-n} \quad \forall n \in \eraum = \nz$} \\
    Was wäre ein fairer Einsatz für dieses Spiel?
    (Wir nennen den Einsatz fair, wenn der Einsatz gleich der erwarteten Auszahlung ist.)\\
    Wir suchen also $\E{\zvx }$

    \txo{$\zvx \geq 0$, EW existiert genau dann, wenn Rechnung endliches Ergebnis hat}

    Nun ist jedoch
    \[
        \sum_{n=1}^{\infty} 2^n \pe(\zvx =2^n)= \sum_{n=1}^{\infty} 2^n \wfk_{\zvx}(2^n)= \sum_{n=1}^{\infty} 2^n 2^{-n} = \sum_{n=1}^{\infty} 1 = \infty
    \]
    Der Erwartungswert existiert also nicht.
    Ein fairer Einsatz kann nicht bestimmt werden.
\end{bsp}


\begin{er} \txc{(2026)} \\
    Aus Satz (E), d.h.\ der Linearität des Erwartungswerts, folgen: % TODO Satz E verlinken
    \[
        \E{a \zvx}= a \E{\zvx } \text{ und } \E{\zvx  + \zvy } = \E{\zvx } + \E{\zvy}
    \]
    für $a \in \rz$ und $\rz$-wertige $\zvx$ und $\zvy$ mit existierendem Erwartungswert
\end{er}

\begin{fr} \txc{(2026)}
Gilt auch \(\E{\zvx \zvy}  = \E{\zvx }  \E{\zvy}\)? \\
    \txo{Spezialfall $\zvx=\zvy$: $\E{E^2} = (\E{\zvx })^2$. \\ Gilt das?
    }
\end{fr}
\begin{bsp} \txc{(2026)}
Es gibt Zufallsvariablen \(\zvx ,\zvy\)mit
    \[
        \E{\zvx \cdot \zvy} \neq \E{\zvx }\cdot\E{\zvy}.
    \]

    Wähle
    \[
        \pe(\zvx =+1)=\pe(\zvx =-1)=\frac{1}{2}
    \]
    und
    \[
        \zvy=\zvx.
    \]

    Dann gilt
    \[
        \E{\zvx }
        =
        \E{\zvy}
        =
        0
    \]
    und
    \[
        \E{\zvx \zvy}
        =
        \E{\zvx ^2}
        =
        \E{1}
        =
        1.
    \]
\end{bsp}

\begin{satz}
    \textbf{Satz (P)} \txc{(2026)}  \\
    Seien
    \[
        \zvx:\eraum\to\mathbb{R}
        \quad\text{und}\quad
        \zvy:\eraum\to\mathbb{R}
    \]
    zwei unabhängige Zufallsvariablen, deren Erwartungswerte existieren.

    Dann existiert auch der Erwartungswert von \(\zvx \cdot \zvy\), und es gilt
    \[
        \E{\zvx \cdot \zvy}
        =
        \E{\zvx }\cdot\E{\zvy}.
    \]

    \paragraph{Beweis.}

    \paragraph{Schritt 1:}
    Wir nehmen zunächst an, dass der Erwartungswert von \(\zvx \cdot \zvy\)existiert.

    Wir beginnen mit der folgenden Darstellung.
    Für \(z\neq 0\)gilt: \txo{(xy = z)}
    \[
        \begin{aligned}
            \pe(\zvx \zvy =z)
            & \overset{\txo{*}}{=}
            \sum_{y\neq 0}
            \pe(\zvx \zvy =z\mid \zvy =y)\,\pe(\zvy=y)
            \\
            &=
            \sum_{y\neq 0}
            P\left(\zvx =\frac{z}{y}\mid \zvy =y\right)
            \pe(\zvy=y)
            \\
            &=
            \sum_{y\neq 0}
            P\left(\zvx =\frac{z}{y}\right)
            \pe(\zvy=y),
        \end{aligned}
    \]
    da \(\zvx \)und \(\zvy\)unabhängig sind.
    \[
        \txo{
            \pe(\zvx  \cdot \zvx = z)= \sum_{y \in \zvy(\eraum)} \pe(\zvx  \cdot \zvy = z, \zvy = y)\overset{z \neq 0}{=} \sum_{y \in \zvy(\eraum), y \neq 0} \pe(\zvx  \cdot \zvy = z| \zvy = y)\cdot \pe(\zvy=y)
        }
    \]
    Aus
    \[
        \pe(\zvx \zvy =z)
        =
        \sum_{y\neq 0}
        P\left(\zvx =\frac{z}{y}\right)
        \pe(\zvy=y),
        \qquad z\neq 0,
    \]
    folgt
    \begin{align*}
        \E{\zvx \zvy}
        &=
        \sum_{z\in \zvx\zvy(\eraum)\txo{, z \neq 0}}
        z\,\pe(\zvx \zvy =z)
        \\
        &=
        \sum_{z\neq 0}
        \sum_{y\neq 0}
        \frac{z}{y}\,y\,
        P\left(\zvx =\frac{z}{y}\right)
        \pe(\zvy=y)\\
        &=
        \sum_{y\neq 0}
        \sum_{x\neq 0}
        xy \pe(\zvx =x)\pe(\zvy=y)\\
        &=
        \left(
                \sum_{y\neq 0} y \pe(\zvy=y)
                \right)
        \left(
                \sum_{x\neq 0} x \pe(\zvx =x)
                \right)\\
        &=
        \E{\zvx } \cdot \E{\zvy}
    \end{align*}

    \paragraph{Schritt 2:} Eine analoge Rechnung zeigt die Existenz des Erwartungswerts:
\end{satz}
\begin{satz}
    \textbf{Kurzfassung Satz (P)} \txc{(2026)}
    \[
        \zvx,\zvy \text{ unabhängig mit ex. EW}
        \;\Longrightarrow\;
        \E{\zvx \zvy}
        =
        \E{\zvx }\,\E{\zvy}.
    \]
\end{satz}

\begin{bsp} \txc{(2026)}
\[
        \E{\zvx \cdot \zvy}
        =
        \E{\zvx }\cdot\E{\zvy}
    \]
    impliziert nicht, dass \(\zvx \)und \(\zvy\)unabhängig sind.

    \begin{itemize}
        \item Definiere Hilfszufallsvariablen \(\zvx _1\)und \(\zvy_1\), unabhängig mit
        \[
            \pe(\zvx _1=0)=\pe(\zvx _1=1)=\frac12
        \]
        und
        \[
            \pe(\zvy_1=0)=\pe(\zvy_1=1)=\frac12.
        \]

        \item Setze
        \[
            \zvx\coloneqq \zvx_1+\zvy_1,
            \qquad
            \zvy\coloneqq|\zvx_1-\zvy_1|.
        \]
    \end{itemize}

    \paragraph{Beobachtungen}

    \begin{itemize}
        \item Bildbereiche:
        \[
            \zvx(\eraum)=\{0,1,2\},
            \qquad
            \zvy(\eraum)=\{0,1\}.
        \]

        \item
        \[
            \zvy =0
            \iff
            \zvx_1=\zvy_1
            \iff
            \zvx\in\{0,2\}.
        \]

        \item
        \[
            \zvy =1
            \iff
            \zvx_1\neq \zvy_1
            \iff
            \zvx =1.
        \]
    \end{itemize}

    \paragraph{\(\zvx \)und \(\zvy\)sind nicht unabhängig.}

    \[
        \pe(\zvx =1,\zvy =0)=0,
    \]
    da
    \[
        \{\zvx =1,\zvy =0\}=\varnothing.
    \]

    Andererseits gilt
    \[
        \pe(\zvx =1)
        =
        2\cdot\frac14
        >
        0
    \]
    und
    \[
        \pe(\zvy=0)
        =
        2\cdot\frac14
        >
        0.
    \]

    \paragraph{Dennoch gilt \(\E{\zvx \zvy}=\E{\zvx }\E{\zvy}\).}

    \[
        \begin{aligned}
            \E{\zvx \zvy}
            &=
            0\cdot \pe(\zvx \zvy =0)
            +
            1\cdot \pe(\zvx \zvy =1)
            +
            2\cdot \pe(\zvx \zvy =2)
            \\
            &=
            0+\pe(\zvx =1)+0
            \\
            &=
            \frac12.
        \end{aligned}
    \]

    Außerdem
    \[
        \begin{aligned}
            \E{\zvx }\,\E{\zvy}
            &=
            (\E{\zvx _1}+\E{\zvy_1})\,\pe(\zvy=1)
            \\
            &=
            2\,\E{\zvx _1}\,\pe(\zvx =1)
            \\
            &=
            2\cdot\frac12\cdot\frac12
            \\
            &=
            \frac12.
        \end{aligned}
    \]
\end{bsp}

\begin{bem} \txc{(2026)} \\
Der Satz kann nicht nur zum Berechnen von Erwartungswerten von Produkten
    unabhängiger Zufallsvariablen verwendet werden, sondern kann auch benutzt
    werden, um zu zeigen, dass zwei Zufallsvariablen nicht unabhängig sein können,
    wenn der Erwartungswert des Produkts der Zufallsvariablen nicht ins Produkt der
    einzelnen Erwartungswerte zerfällt.
\end{bem}


\begin{defi}
    \textbf{Zusammengesetzte Zufallsvariablen} \txc{(2026)}  \\
    Eine Zufallsvariable $\zvx: \eraum \rightarrow \rz$ und eine \txo{messbare} Funktion $f: \rz \rightarrow \rz$ können
    zusammengesetzt werden zu einer neuen Zufallsvariablen:
    \[
        \zvy(\ele)= (f \circ  \zvx)(\ele)= f(\zvx (\ele))
    \]
    \txo{
        $\zvy(\ele)= f(\zvx (\ele))\quad \forall \ele$
    }
\end{defi}

\begin{satz}
    \label{satz:7-4} \txc{(2026)} \\
    Sofern der Erwartungswert $\E{\zvy}$ existiert, gilt
    \[
        \E{\zvy} = \sum_{x \in \zvx(\eraum)} f(x)\pe(\zvx =x)
    \]
    \txo{
        bzw., falls $\zvx$ Dichte d hat und EW von $\zvy$ gleich $f(x)$ existiert
    }
    \[
        \txo{
            \E{\zvy} = \intinfs f(x)d(x)\dx
        }
    \]

    \paragraph{Beweis.}
    \begin{align*}
        \E{\zvy} & =   \sum_{y \in \zvy(\eraum)} y \pe(\zvy=y)\\
        & =   \sum_{y \in \zvy(\eraum)} y  \sum_{x \in \zvx(\eraum): f(x)= y\txo{*}} \pe(\zvx =x)\\
        & =   \sum_{y \in \zvy(\eraum)}   \sum_{x \in \zvx(\eraum): f(x)= y} \oub{y}{=f(x)} \pe(\zvx =x)\\
        & =\sum_{y \in \zvy(\eraum)}   \sum_{x \in \zvx(\eraum): f(x)= y} f(x)\pe(\zvx =x)\\
        & = \sum_{x \in \zvx(\eraum)} f(x)\pe(\zvx =x)\\
        & \txo{= \E{f(x)}}
    \end{align*}
    \[
        \txo{*  \zvy =y \Leftrightarrow f(x)= y, \quad \pe(\zvy=y)= \sum_{x \in \zvx(\eraum): f(x)= y} \pe(\zvx =x)}
    \]
\end{satz}

\begin{bsp}
    Zu Satz~\ref{satz:7-4}. \\
    Mit $f(x)= x^2$
\end{bsp}

\karkar{satz_linearitaet_erwartungswert}
\karkar{bem_linearitaet_erspart_direkte_rechnung}
\karkar{def_zusammengesetzte_zufallsvariable}
\karkar{satz_transformationssatz_erwartungswert}
\karkar{frage_gilt_ew_produkt_gleich_produkt_ew}
\karkar{satz_p_erwartungswert_produkt_unabhaengig}
\karkar{bem_satz_p_zum_widerlegen_unabhaengigkeit}
\karkar{bem_umkehrung_satz_p_gilt_nicht}
\karkar{bem_ew_muss_nicht_existieren}


%%%%%%%%%%%%%%%%%%%%%%%
% Varianz
%%%%%%%%%%%%%%%%%%%%%%%

\section{Varianz}

\begin{defi}
    \textbf{n-tes Moment von $\zvx$} \txc{(2026)}  \\
    Für $n \in \nz $ heißt $\E{\zvx ^n}$ das erste n-te Moment von $\zvx$, insofern $\E{|\zvx|^n} < \infty$.
    (Erwartungswert ist das 1.
    Moment)
\end{defi}
\begin{defi} \txc{(2026)}
$\E{ (\zvx -  \E{\zvx })^n } $ heißt das n-te zentrierte Moment von $\zvx$.
Es existiert, wenn $\E{|\zvx|^n} < \infty$. \\
\end{defi}

\begin{defi} \txc{(2026)}
Die Varianz von $\zvx$ ist das 2 zentrierte Moment, d.h.
    \[
        Var(\zvx)= \E{ (\zvx -  \E{\zvx })^2 }
    \]
    \txo{(mittlere quadratische Abweichung)} \\
    Die Varianz von $\zvx$ existiert, wenn $Var(\zvx)< \infty$. \\
    Die Varianz einer Zufallsvariable $\zvx$ mit existierendem Erwartungswert ist gegeben
    durch
    \[
        Var(\zvx)= \E{ (\zvx -  \E{\zvx })^2 } = \sum_{x\in \zvx(\eraum)} (x - \E{\zvx })^2 \pe(\zvx =x)
    \]
    (gewichtete quadratische Abweichung vom Erwartungswert)
    \txo{bzw.}
    \[
        \txo{\Var(\zvx)= \intinfs \underbrace{ (x - \E{\zvx })^2 }{=f(x)} d(x)\dx }
    \]
\end{defi}

\begin{bem} \txc{(2026)}
$\Var(\zvx)\geq 0$
\end{bem}

\begin{defi}
    \textbf{Existierende Varianz} \txc{(2025)}  \\
    Wir sagen, dass die Varianz existiert, wenn
    \[
        \Var(\zvx)= \sum_{x\in \zvx(\eraum)} (x - \E{\zvx })^2 \pe(\zvx =x)< \infty
    \]
    \txo{bzw.}
    \[
        \txo{ \intinfs (x - \E{\zvx })^2 d(x)\dx < \infty }
    \]
\end{defi}

\begin{bsp}
    \textbf{Bsp.} \txc{(2026)} \\

    Seien \(\zvx , \zvy\)Zufallsvariablen mit
    \[
        \pe(\zvx =+1)=\frac{1}{2}=\pe(\zvx =-1)
    \]
    und
    \[
        \zvy =100\cdot \zvx.
    \]

    \textbf{Anhand des Erwartungswertes: kein Unterschied!}

    \[
        \E{\zvx }
        =1\cdot \pe(\zvx =1)+(-1)\cdot \pe(\zvx =-1)
        =1\cdot\frac{1}{2}-1\cdot\frac{1}{2}
        =0.
    \]

    Da \(\zvy=100 \zvx\)gilt,

    \[
        \E{\zvy}=100\cdot \E{\zvx }=0.
    \]
    \textbf{Die Varianz zeigt den Unterschied:}

    da \(\E{\zvx }=0\),

    \[
         \Var(\zvx)
        =\E{(\zvx - \oub{\E{\zvx }}{=0})^{2}}
        =\E{ \oub{ \zvx^{2}}{1} }
        = \E{1}
        = 1
    \]

    Wegen \(\zvx \in\{-1,1\}\)ist

    \[
         \zvx^{2}=1,
    \]

    also

    \[
         \Var(\zvx)
        =\E{1}
        =1.
    \]

    Ebenso gilt wegen \(\E{\zvy}=0\),

    \[
         \Var(\zvy)
        =\E{(\zvy-\E{\zvy})^{2}}
        =\E{\zvy^2}
        =\E{(100 \zvx)^{2}}
    \]

    Damit folgt

    \[
         \Var(\zvy)
        =\E{100^{2} \zvx^{2}}
        =100^{2}\E{\zvx ^{2}}
        =100^{2}.
    \]
\end{bsp}

\begin{bsp}
    \textbf{Bernoulli-Variable (Spezialfall der Binomialverteilung mit n = 1)} \txc{(2026)}  \\
    Zufallsvariable $\zvx$ mit $\pe(\zvx  = +1)= p = 1 - \pe(\zvx =0)$ für Erfolgswahrscheinlichkeit $p \in [0,1]$ und $\E{\zvx } = p$
    . \\
    \txc{$\Var(\zvx)$ existiert, da}
    \begin{align*}
        \txo{\Var(\zvx)} & \txo{= \sum_{x\in \zvx(\eraum)}  (x - \E{\zvx })^2 \pe(\zvx =x)} \\
        &
        \txo{ = (0 - \underbrace{\E{\zvx }}{=p})^2  \underbrace{\pe(\zvx =0)}{=1-p} + (1 - \underbrace{\E{\zvx }}{=p})^2  \underbrace{\pe(\zvx =1)}{=p}} \\
        & \txo{ = p^2 (1-p)+ (1-p)^2 p } \\
        & \txo{= p(1-p)[p+(1-p)]} \\
        & \txo{= p(1-p)< \infty}
    \end{align*}
    \begin{align*}
        \Var(\zvx)
        & = (1 - \E{\zvx })^2 p + (0 - \E{\zvx })^2 (1 - p)\\
        & = (1 - p)^2 p + p^2 (1 - p)\\
        & = p(1 - p)
    \end{align*}
\end{bsp}

\begin{bsp}
    \textbf{gleichverteilte Zufallsvariable auf \{1,..,n\}} \txc{(2026)}  \\
    Zufallsvariable $\zvx$ mit $\pe(\zvx  = k)= \frac{1}{n} \text{ für } k = 1, \dots, n$

    \[
        \E{\zvx }
        = \sum_{k=1}^{n} k \cdot \frac{1}{n}
        = \frac{n(n+1)}{2n}
        = \frac{n+1}{2}
    \]
    und
    \[
        \Var(\zvx)= \frac{n^2 - 1}{12}
    \]
    weil
    \begin{align*}
        \Var(\zvx)
        &= \frac{1}{n} \sum_{k=1}^{n} (k - \E{\zvx })^2 \\
        &= \frac{1}{n} \sum_{k=1}^{n} \left(k^2 - 2k\E{\zvx } + \E{\zvx }^2\right)\\
        &= \frac{1}{n} \sum_{k=1}^{n} k^2 - 2\E{\zvx } \cdot \oub{\frac{1}{n} \sum_{k=1}^{n} k}{=\E{\zvx }} + \E{\zvx }^2 \\
        &= \frac{(n+1)(2n+1)}{6} - \E{\zvx }^2 \\
        &= \frac{n^2 - 1}{12}
    \end{align*}
\end{bsp}

\begin{lem}
    \textbf{Technisches Lemma} \txc{(2026)}  \\
    Für alle $a \in \rz$ gilt $\E{(\zvx  - a)^2} = \Var(\zvx)+ (\E{\zvx } - a)^2$

    \paragraph{Beweis} TODO F 76 % TODO
\end{lem}

\begin{folg}
    \textbf{Folgerung I} \txc{(2026)}  \\
    $a \Leftrightarrow \E{(\zvx  - a)^2}$ wird minimal für $a = \E{\zvx }$.

    In diesem Fall ist $\E{(\zvx  - a)^2} = \Var(\zvx)$ per Definition der Varianz.
\end{folg}

\begin{folg}
    \textbf{Folgerung II und wichtige Rechenregel} \txc{(2026)}  \\
    Die Wahl a = 0 zeigt
    \[
        \txo{0 \leq } \Var(\zvx)= \E{\zvx ^2} - (\E{\zvx })^2
    \]
    Es gilt daher stets $(\E{\zvx })^2 \le \E{\zvx ^2}$.
\end{folg}

\karkar{def_n_tes_moment}
\karkar{def_n_tes_zentriertes_moment}
\karkar{def_varianz}
\karkar{bem_varianz_nichtnegativ}
\karkar{def_existenz_varianz}
\karkar{lem_technisches_lemma_ew_verschobenes_quadrat}
\karkar{folg_minimierungseigenschaft_varianz}
\karkar{folg_verschiebungssatz_varianz}

%%%%%%%%%%%%%%%%%%%%%%%
% Rechenregeln für die Varianz
%%%%%%%%%%%%%%%%%%%%%%%

\subsection{Rechenregeln für die Varianz}

\begin{satz}
    \textbf{Satz (V)} \label{satz:satz-v} \txc{(2026)}  \\
    Seien $\zvx, \zvy$ diskrete\txc{(?)})Zufallsvariablen mit existierendem Erwartungswert, $a,b \in \rz$.

    \paragraph{(a)} $\Var(a \zvx+b)= a^2 \Var(\zvx)$

    \begin{align*}
        \txo{\Var(a \zvx+b)} & \txo{ \E{ (a \zvx+b - \E{a \zvx+b})^2 } } \\
        & \txo{ \E{ (a (\zvx -\E{\zvx }))^2 } } \\
        & \txo{ a^2 \E{(\zvx - \E{\zvx })^2} } \\
        & \txo{= a^2 \Var(\zvx)}
    \end{align*}


    \paragraph{(b)} $\Var(\zvx  + \zvy)= \Var(\zvx)+ \Var(\zvy)+ 2 \Cov(\zvx  , \zvy)$ mit Kovarianz
    $\Cov(\zvx ,\zvy)\coloneqq \E{(\zvx -\E{\zvx })(\zvy-\E{\zvy})}$.
    \begin{align*}
        \txo{\Var(\zvx  + \zvy)} & \txo{= \E{ ((\zvx +\zvy)- \E{(\zvx +\zvy)})^2 } } \\
        & \txo{= \E{ (\zvx -\E{\zvx })^2 }  + 2 \E{ (\zvx -\E{\zvx })(\zvy-\E{\zvy})} + \E{(\zvy-\E{\zvy})^2}  } \\
        & \txo{  \Var(\zvx)+ \Var(\zvy)+ 2 \Cov(\zvx  , \zvy)}
    \end{align*}

    \paragraph{(c)} $\zvx$ und $\zvy$ unabhängig $\Rightarrow \Var(\zvx  + \zvy)= \Var(\zvx)+ \Var(\zvy)$.
    \begin{align*}
        \txo{\text{$\zvx$ und $\zvy$ unabhängig}} & \txo{\Rightarrow \zvx-\E{\zvx }, \zvy-\E{\zvy}\ unabh.} \\
        &  \txo{\Cov(\zvx ,\zvy)} & \txo{= \E{(\zvx -\E{\zvx })(\zvy-\E{\zvy})}} \\
        & & \txo{ \overset{(P)}{=}  \E{\zvx -\E{\zvx }} \E{\zvy-\E{\zvy}}} \\
        & & \txo{=0}
    \end{align*}

    \paragraph{(d)} Für diskrete Zufallsvariablen gilt: \\
    $\Var(\zvx)= 0 \Rightarrow$ Es existiert eine Konstante c mit $\pe(\zvx =c)= 1$. \\
    $c = \E{\zvx }$

    \paragraph{Beweis für (a)bis (c):} Verwende $\Var(Z)= \E{(Z- \E{Z})^2}$ und die Rechenregeln für den EW\@.

    \paragraph{Beweis für (d):} $0 = \Var(\zvx)= \E{(\zvx - \E{\zvx })^2} = \sum_{x\in \zvx(\eraum)} (x- \E{\zvx })^2 \pe(\zvx =x)$. \\
    Alle Summanden sind $\geq 0$, folglich impliziert $\pe(\zvx =x)> 0$ sofort $x = \E{\zvx }$
    \txc{, da die Summe nicht 0 sein könnte, wenn ein Summand nicht 0 wäre.
    Also muss jedes $(x- \E{\zvx })^2 \pe(\zvx =x)$ sein und wenn $\pe(\zvx =x)> 0$, muss also $(x- \E{\zvx })^2 = 0$ sein, also muss $x- \E{\zvx } = 0$ sein, also $x = \E{\zvx }$}.

    Damit ist gezeigt, dass $\pe(\zvx =\E{\zvx })= 1$\txc{, da jedes $x \in \zvx$ gleich $\E{\zvx }$ ist}.

    \txo{Kann nicht gelten, wenn f Dichte hat}
\end{satz}

\karkar{def_kovarianz}
\karkar{satz_v_a_varianz_affine_transformation}
\karkar{satz_v_b_varianz_summe_allgemein}
\karkar{satz_v_c_varianz_summe_unabhaengig}
\karkar{satz_v_d_varianz_null_konstante}
% TODO F 82 Was sol ldas mit dem herumkritzeln?
% TODO Auf welchen regeln diese Übergänge basieren

%%%%%%%%%%%%%%%%%%%%%%%
% Varianzen der wichtigsten Verteilungen
%%%%%%%%%%%%%%%%%%%%%%%


%\subsection{Varianzen der wichtigsten Verteilungen}

\begin{bsp}
    \textbf{Binomialverteilte Zufallsvariable $S_n$ mit Parametern n und p} \txc{(2026)}  \\
    $\eraum = \{0,1\}^n$ und

    \[
        p(\ele)
        = p^{\sum_{i=1}^{n} \ele_i}
        (1-p)^{\,n-\sum_{i=1}^{n} \ele_i}
    \]
    Darstellung:
    \[
        S_n = \sum_{i=1}^{n} \zvx_i
    \]
    mit
    \[
        \zvx_i(\ele)= \ele_i
        \quad \text{für alle } \ele \in \eraum.
    \]

    Dabei sind die $\zvx_i$ unabhängig mit
    \[
        \pe(\zvx _i = 1)= p
    \]
    und
    \[
        \pe(\zvx _i = 0)= 1-p = q
    \]
    für alle $i$.

    \[
        \E{\zvx }
        = \sum_{i=1}^{n} \E{\zvx _i}
        = np
    \]

    nach Satz (E). % TODO Satz E verlinken
\end{bsp}

Unter Verwendung der Unabhängigkeit der $\zvx_i$ folgt ohne weitere Arbeit aus Satz (V)(\ref{satz: satz-v})
\[
    \Var(\zvx)= \sum^n_{i=1} \oub{\Var(\zvx _i)}{=pq} = npq
\]

\begin{bsp}
    \textbf{Poisson-verteilte Zufallsvariable $\zvx$ mit Parameter $\lambda > 0$} \txc{(2026)}
    \[
        \E{\zvx } = \lambda = \Var(\zvx)
    \]
    Trick:
    \[
        \Var(\zvx)= \E{\zvx ^2} - (\E{\zvx })^2 = \E{\zvx (\zvx -1)} + \E{\zvx } - (\E{\zvx })^2
    \]
    Wir wissen bereits:
    \[
        \E{\zvx } = \lambda \quad und \quad  (\E{\zvx })^2 = \lambda^2
    \]
    Die Behauptung folgt daher aus
    \[
        \E{\zvx (\zvx -1)}
        = \sum_{k=2}^{\infty} k(k-1)\,\frac{\lambda^k}{k!} e^{-\lambda}
        = \lambda^2 \sum_{l=0}^{\infty} \frac{\lambda^l}{l!} e^{-\lambda}
        = \lambda^2
    \]

    und damit
    \[
        \Var(\zvx)
        = \E{\zvx (\zvx -1)} + \E{\zvx } - (\E{\zvx })^2
        = \lambda^2 + \lambda - \lambda^2
        = \lambda
    \]
\end{bsp}

\begin{bsp}
    \textbf{geometrisch verteilte Zufallsvariable $\zvx$ mit Erfolgswahrscheinlichkeit p} \txc{(2026)}  \\
    \[
        \E{\zvx } = \frac{1}{p} \quad und \quad \Var(\zvx)= \frac{1-p}{p^2}
    \]
    \txo{(Nicht auswendig lernen)}
\end{bsp}

\begin{bsp}
    \textbf{auf [a, b] uniform verteilte Zufallsvariable $\zvx$} \txc{(2026)}  \\
    Dann ist
    \[
        \E{\zvx } = \frac{a + b}{2}
    \]
    und
    \begin{align*}
        \E{\zvx ^2}&  = \intab x^2 \frac{1}{b - a} \dx \\
        &= \frac{1}{3(b - a)} \ x^3 \Big|_{x=a}^b \\
        &= \frac{b^3 - a^3}{3(b - a)} \\
        &= \frac{1}{3}(a^2 + ab + b^2)
    \end{align*}
    \[
        \Var(\zvx)= \E{\zvx ^2} - (\E{\zvx })^2
        = \frac{a^2 + ab + b^2}{3} - \frac{a^2 + 2ab + b^2}{4}
        = \frac{(a - b)^2}{12}^{\txo{*}}
    \]
    \txo{* Nicht auswendig lernen.}
\end{bsp}

\begin{bsp}
    \textbf{Exponentialverteilte Zufallsvariable $\zvx$} \txc{(2026)}
    \[
        \E{\zvx } = \frac{1}{\alpha}
    \]
    \begin{align*}
        \E{\zvx ^2}  & = \int^{\infty}_0 \underbrace{x^2}_{\downarrow} \underbrace{\alpha e^{- \alpha x}}_{\uparrow} \dx \\
        & = - \lim_{N \rightarrow \infty} x^2 e^{- \alpha x} \Big|_{x=0}^N + 2 \int_0^{\infty} x e^{- \alpha x} \dx \\
        & = 0 + \frac{2}{\alpha} \E{\zvx } \\
        & = \frac{2}{\alpha^2}
    \end{align*}
    \[
        \Var(\zvx)= \E{\zvx ^2} - \E{\zvx }^2 =  \frac{2}{\alpha^2} - \left(\frac{1}{\alpha} \right)^2 =  \frac{1}{\alpha^2}
    \]
\end{bsp}

\untit{Ende der Vorlesung vom 5. Juni 2026}

%%%%%%%%%%%%%%%%%%%%%%%%%%%%
% Dichte von  \zvx^2
%%%%%%%%%%%%%%%%%%%%%%%%%%%%

\begin{satz} \txc{(2026)}
    Sei $f_X$ eine Dichte der Zufallsvariablen $\zvx$ . Dann gelten:
\begin{itemize}
    \item[(a)] Die Zufallsvariable $\zvx^2$ hat die Dichte
    \[
        f_{ \zvx^2} (y)= \begin{cases}
                          \frac{1}{2 \sqrt{y}} \left[ f_X \left(\sqrt{y}\right)+ f_X  \left(-\sqrt{y}\right)\right] & y > 0 \\
                          0 & sonst
        \end{cases}
    \]
    \item[(b)] \(\E{\zvx ^2} = \intinfs x^2 f_{\zvx}(x)\dx  \)
\end{itemize}
\paragraph{Beweis (a).}
    Für $y \leq 0$ $F_{ \zvx^2} (y)= \pe (\zvx ^2 \leq y)= 0$. \\
    Sei nun $y > 0$.
    Die Substitution $x = \sqrt{z}$ zeigt:
    \begin{align*}
        F_{ \zvx^2}(y)& = \pe (\zvx ^2 \leq y)\\
        & = \pe (- \sqrt{y} < \zvx \leq \sqrt{y})\\
        & = \int_0^{\sqrt{y}} f_{\zvx}(x)\dx + \int^0_{-\sqrt{y}} f_{\zvx}(x)\dx \\
        & = \int_0^{\sqrt{y}} f_{\zvx}(x)\dx + \int_0^{\sqrt{y}} f_{\zvx}(- x)\dx \\
        \int_0^{y} \frac{1}{2 \sqrt{z}} \left[  f_X (\sqrt{z})+ f_X (-\sqrt{z})\right] dz
    \end{align*}
\paragraph{Beweis (b).}
Die Definition des Erwartungswerts und die Substitutionen $y = x ^2$ sowie
$z = -x$, die die Substitutionen in Teil (a)rückgängig machen, zeigen
\begin{align*}
    \E{\zvx ^2} & = \intinfs y f_{ \zvx^2} (y)dy \\
    & = \int_{0}^{\infty} \frac{y}{2\sqrt{y}}  \left[  f_X (\sqrt{y})+ f_X (-\sqrt{y})\right] dy \\
    & = \int_{0}^{\infty} \frac{x}{2} \left[  f_X (\sqrt{y})+ f_X (-\sqrt{y})\right] 2x \dx  \\
    & = \int_{0}^{\infty} x^2 f_X (x)\dx + \int^0_{-\infty}  z^2 f_{\zvx}(z)dz \\
    & = \intinfs x^2 f_X (x)\dx
\end{align*}
\end{satz}

\karkar{satz_dichte_von_x_quadrat}
\karkar{bem_ex2_direkt_ueber_dichte_von_x}

\section*{Karteikarten Erwartungswert und Varianz der wichtigsten Verteilungen}
\karkar{bem_ew_var_binomialverteilung}
\karkar{bem_ew_var_bernoulli}
\karkar{def_geometrische_verteilung_zaehldichte}
\karkar{bem_ew_var_geometrische_verteilung}
\karkar{def_poisson_verteilung}
\karkar{bem_ew_var_poisson_verteilung}
\karkar{bem_ew_var_uniform_diskret}
\karkar{bem_ew_var_uniform_stetig}
\karkar{def_exponentialverteilung}
\karkar{bem_ew_var_exponentialverteilung}
\karkar{bem_technik_binomial_ew_var_via_indikatoren}
\karkar{bem_technik_faktorielle_momente_varianz}


% Ab F 101 ist es einfach eine versehentliche Wdh der früheren Folien
